\documentclass[a4paper,12pt,titlepage]{article} %page properties
\usepackage[headings]{fullpage} %for fullpage with headings
\usepackage{fancyhdr} %for headings
\usepackage{graphicx} %for label
\usepackage{tabularx} %for fullpage tables
\usepackage{longtable} %for tables on multiple pages

\usepackage{hyperref}

% defining headings
\pagestyle{fancy}
\fancyhead{}
\fancyhead[LE,RO]{PAMI 2015-2016 \hfill HOMEWORK REPORT}

\begin{document}
\begin{titlepage}

%defining university
\begin{center}
	POLITECNICO DI MILANO --- COMO CAMPUS\\
	\vspace{10pt}
	\includegraphics[scale=0.1]{logo-polimi.png}\\
	\vspace{10pt}
	PATTERN ANALYSIS AND MACHINE INTELLIGENCE 2015-2016\\
	prof. Matteo Matteucci	
	\line(2,0){500}
\end{center}

\vspace{60pt}

%defining work type
\begin{center}
	{\Huge \textbf{Homework report}}\\
\end{center}

\vspace{60pt}

%defining project repository
\begin{center}
	{\large Project Repository}
\end{center}
\begin{tabularx}{\textwidth}{|X|}
	\hline
	\href{https://github.com/attillax/PAMI-2015}{Click}\\
	\hline
\end{tabularx}

\vspace{20pt}

%defining team members
\begin{center}
	{\large Team Members}
\end{center}
\begin{tabularx}{\textwidth}{|X|X|X|}
	\hline
	ID & Surname & Name\\
	\hline
	10460625 & Golubeva & Svetlana\\
	\hline
	 &  & \\
	\hline
	 &  & \\
	\hline
	 &  & \\
	\hline
\end{tabularx}

\vspace{\fill}
\begin{center}
	\line(2,0){500}
\end{center}

\end{titlepage}

% % % % % % % % % % % % % % % % % % % % % % % % % % % % % % % % %
\tableofcontents

% % % % % % % % % % % % % % % % % % % % % % % % % % % % % % % % %
\newpage
\listoftables

\listoffigures

% % % % % % % % % % % % % % % % % % % % % % % % % % % % % % % % %
\newpage
\section{Homework}

% % % % % % % % % % % % % % % % % % % % % % % % % % % % % % % % %
\newpage
\section{Technical Notes}

\subsection{Chronology}
\begin{longtable}{|c|p{13cm}|}
	\hline
	Date & Content\\
	\hline
	&  Release\\
	\hline
	&  Add to repoitory\\
	\hline
	&  Solution\\
	\hline
	&  Submission\\
	\hline
	&  \\
	\hline
	& \\
	\hline
	& \\
	\hline
	\caption{project's chronology.}
\end{longtable}

\subsection{List of tools}
\begin{table}[h]
	\begin{tabularx}{\textwidth}{|l|X|}
		\hline
		What & Which \\
		\hline
		OS & Linux \\
		\hline
		Lang & R, \TeX \\
		\hline
		IDE & R-Studio, \TeX-Studio, Emacs\\
		\hline
		& Terminal\\
		\hline
		& \\
		\hline
	\end{tabularx}
	\caption{list of tools.}
\end{table}

% % % % % % % % % % % % % % % % % % % % % % % % % % % % % % % % %
\end{document}
